% spielkarten-inhalt.tex — die Karten selbst, ohne Praeambel.
%
% Eingebunden von beispiel/spielkarten.tex (Einzelkarten) und von
% beispiel/spielkarten-bogen.tex (A4-Druckbogen). Beide setzen denselben
% Satz Karten; sie unterscheiden sich nur in der Klassenoption. Genau
% dafuer liegt der Inhalt in einer eigenen Datei — so ist am gebauten PDF
% nachzupruefen, dass die Druckfassung am Satz der Karte nichts aendert.
%
% Zwoelf Karten, jeder der acht Kartentypen mindestens einmal, in vier
% Randfarben. Zwoelf und nicht neun: so laeuft der Druckbogen auf ein
% zweites Bogenpaar ueber, und dass das stimmt, ist am PDF zu sehen.
%
% Der Inhalt ist erfunden. „Die Uferlande" gibt es in Aventurien nicht;
% Namen, Werte und Regeln sind frei erdacht und stehen hier nur, damit
% jedes Element einen Text hat, der so lang ist wie im Ernstfall.
%
% Die Abbildungen sind Platzhalter und werden erzeugt mit
%
%   python3 werkzeuge/kartenplatzhalter.py "Bruder Halmrich" "Zwergenspalter" \
%       "Schuppenpanzer" "Wegfeuerlaterne" "Moorschrecke" "Windhauch" \
%       "Klingentanz" "Uferlande" "Gerlinde Rabenfels" "Nebelschleicher" \
%       "Wurzelsegen" "Steinwurf"
%   python3 werkzeuge/kartenplatzhalter.py --rund "Halmrich" "Gerlinde" \
%       "Moorschrecke" "Nebelschleicher"
%
% Die Randfarben kommen aus
%
%   python3 werkzeuge/kartengrafik.py "/pfad/zum/Kartenpaket" \
%       --farbe blau --farbe rot --farbe gruen --farbe bernstein

% ======================================================================
% Erste Gruppe: die Bewohner der Uferlande — blauer Rand
% ======================================================================

\dsaKartenfarbe{blau}

% --- NSC --------------------------------------------------------------

\begin{dsaKarteNSC}
  \dsaKartenname{Bruder Halmrich}
  \dsaKartenuntertitel{Wandernder Efferdgeweihter}
  \dsaKartenbild{platzhalter-bruder-halmrich}
  \dsaKartennummer{UFR 01}
  \dsaKartenmedaillon{platzhalter-halmrich-rund}
  \dsaKartennachweis{Platzhalter}

  \dsaKartentext{Bedächtiger Zuhörer, Götter \& Kulte (12/13/11) 9,
    Boote \& Schiffe (11/12/13) 8}
  \dsaKartentext{Zieht seit dreißig Jahren den Uferpfad entlang und segnet
    die Netze der Fischer. Spricht langsam und wiegt jedes Wort zweimal.
    Hat eine unerklärliche Schwäche für geräucherten Aal und eine ebenso
    unerklärliche Abneigung gegen Brückenzölle.}
  \dsaKartenstich{Charakterzüge}{gelassen, wortkarg, unbestechlich}
  \dsaKartenstich{Dienstleistungen}{Wettervorhersage, Segnung von Booten,
    Vermittlung bei Streit unter Fischern}
  \dsaKartenstich{Im Abenteuer}{Halmrich hat am Ufer eine Truhe gefunden,
    die er nicht öffnen mag. Er sucht jemanden, der sie zum Tempel nach
    Wehrheim bringt — ohne hineinzusehen.}
  \dsaKartenstich{Zusammenfassung}{Mensch, Uferlande, Geweihter, Auftrag}
\end{dsaKarteNSC}

\begin{dsaKarteNSC}
  \dsaKartenname{Gerlinde Rabenfels}
  \dsaKartenuntertitel{Wirtin zum Krummen Ruder}
  \dsaKartenbild{platzhalter-gerlinde-rabenfels}
  \dsaKartennummer{UFR 02}
  \dsaKartenmedaillon{platzhalter-gerlinde-rund}
  \dsaKartennachweis{Platzhalter}

  \dsaKartentext{Schlagfertige Wirtin, Gassenwissen (13/14/12) 11,
    Zechen (12/12/14) 9}
  \dsaKartentext{Führt die einzige Schenke zwischen Wehrheim und der
    Moorkante. Weiß alles, was am Fluss geredet wird, und gibt davon genau
    so viel weiter, wie ihr nützt.}
  \dsaKartenstich{Charakterzüge}{schlagfertig, misstrauisch, großzügig
    gegenüber Kindern}
  \dsaKartenstich{Dienstleistungen}{Unterkunft, Gerüchte, Botendienste
    flussaufwärts}
  \dsaKartenstich{Im Abenteuer}{Seit drei Wochen zahlt ein Gast in alten
    Münzen, die niemand kennt. Gerlinde will wissen, woher sie kommen —
    fragen mag sie aber nicht selbst.}
  \dsaKartenstich{Zusammenfassung}{Mensch, Uferlande, Wirtin, Information}
\end{dsaKarteNSC}

% ======================================================================
% Zweite Gruppe: Ausrüstung — roter Rand
% ======================================================================

\dsaKartenfarbe{rot}

% --- Waffe ------------------------------------------------------------

\begin{dsaKarteWaffe}
  \dsaKartenname{Zwergenspalter aus Wehrheim}
  \dsaKartenbild{platzhalter-zwergenspalter}
  \dsaKartennummer{UFR 03}

  \dsaKartenwert{KT}{Zweihandhiebwaffen}
  \dsaKartenwert{TP}{\dsaFormel{1W6+5}}
  \dsaKartenwert{L+S}{KK 15}
  \dsaKartenwert{AT/PA-Mod}{\dsaFormel{0/\dsaMinus 1}}
  \dsaKartenwert{RW}{mittel}
  \dsaKartenwert{Gewicht}{3,5 Stn}
  \dsaKartenwert{Länge}{110 HF}
  \dsaKartenwert{Preis}{190 S}
  \dsaKartenwert{Komplexität}{komp (1 AP)}

  \dsaKartenstich{Waffenvorteil}{Gegen Rüstungen mit RS~4 oder höher
    richtet der Zwergenspalter \dsaMod{+}{2}~TP an. Die Klinge ist dafür
    eigens gezahnt.}
  \dsaKartenstich{Waffennachteil}{Die Waffe ist kopflastig. Nach einer
    misslungenen Attacke ist die nächste Verteidigung um 1 erschwert.}
\end{dsaKarteWaffe}

% --- Rüstung ----------------------------------------------------------

\begin{dsaKarteRuestung}
  \dsaKartenname{Schuppenpanzer der Flussgarde}
  \dsaKartenbild{platzhalter-schuppenpanzer}
  \dsaKartennummer{UFR 04}

  \dsaKartenwert{Rüstungsschutz}{5}
  \dsaKartenwert{Belastung (Stufe)}{2}
  \dsaKartenwert{zusätzliche Abzüge}{\dsaFormel{\dsaMinus 1} auf GE}
  \dsaKartenwert{Gewicht}{9 Stn}
  \dsaKartenwert{Preis}{340 S}
  \dsaKartenwert{Komplexität}{komp (2 AP)}

  \dsaKartenstich{Rüstungsvorteil}{Die Schuppen sind gefettet. Im Wasser
    gilt die Belastung um eine Stufe verringert, solange der Träger sich
    nicht länger als eine Spielrunde untertaucht.}
  \dsaKartenstich{Rüstungsnachteil}{Gegen Waffen der Kampftechnik
    Stangenwaffen schützt der Panzer nur mit RS~3: die Schuppen geben
    einer geführten Spitze nach.}
\end{dsaKarteRuestung}

% --- Gegenstand -------------------------------------------------------

\begin{dsaKarteGegenstand}
  \dsaKartenname{Wegfeuerlaterne}
  \dsaKartenbild{platzhalter-wegfeuerlaterne}
  \dsaKartennummer{UFR 05}

  \dsaKartenstich{Beschreibung}{eine Laterne aus getriebenem Kupfer, deren
    Scheiben mit Efferdsrunen geätzt sind}
  \dsaKartenstich{Wirkung}{Ihr Licht durchdringt Nebel bis auf zwanzig
    Schritt. Wer die Laterne trägt, erhält eine Erleichterung von 2 auf
    Orientierung, solange er einem Wasserlauf folgt.}
  \dsaKartenstich{Kosten}{drei Wochen Arbeit eines Kupferschmieds}
  \dsaKartenstich{Preis}{280 S}
\end{dsaKarteGegenstand}

\begin{dsaKarteGegenstand}
  \dsaKartenname{Steinwurfnetz}
  \dsaKartenbild{platzhalter-steinwurf}
  \dsaKartennummer{UFR 06}

  \dsaKartenstich{Beschreibung}{ein Wurfnetz mit eingeflochtenen
    Bleiperlen, wie es die Aalfischer benutzen}
  \dsaKartenstich{Wirkung}{Bei einer gelungenen Fernkampfprobe erhält ein
    Ziel der Größenkategorie klein oder mittel den Status
    \emph{Gefesselt}. Befreien gelingt mit einer Probe auf Kraftakt,
    erschwert um 2.}
  \dsaKartenstich{Preis}{24 S}
\end{dsaKarteGegenstand}

% ======================================================================
% Dritte Gruppe: Zauberwerk und Geweihtes — grüner Rand
% ======================================================================

\dsaKartenfarbe{gruen}

% --- Zauber, kurz: beide Seiten tragen denselben Text ------------------

\begin{dsaKarteZauber}
  \dsaKartenname{Windhauch}
  \dsaKartenuntertitel{Zaubertrick}
  \dsaKartennummer{UFR 07}

  \dsaKartentext{Ein kurzer Luftzug fährt über das Wasser und legt sich
    wieder. Er reicht, um ein Segel einmal knattern zu lassen, eine Kerze
    zu löschen oder ein loses Blatt einen Schritt weit zu tragen.}
  \dsaKartenstich{Reichweite}{16 Schritt}
  \dsaKartenstich{Wirkungsdauer}{sofort}
  \dsaKartenstich{Zielkategorie}{Objekte}
  \dsaKartenstich{Merkmal}{Elementar (Luft)}
  \dsaKartenstich{Anmerkung}{Elfen, Hexen}
\end{dsaKarteZauber}

% --- Zauber, lang: mit Fortsetzung auf der Rückseite -------------------

\begin{dsaKarteZauber}
  \dsaKartenname{Wurzelsegen}
  \dsaKartenuntertitel{Liturgie}
  \dsaKartennummer{UFR 08}

  \dsaKartentext{Der Geweihte legt die Hand auf den Boden und ruft die
    Kraft des Flussvaters an. Im Umkreis von QS Schritt treiben Wurzeln
    aus dem Erdreich und halten fest, was sie fassen können.}
  \dsaKartenstich{Liturgiezeit}{4 Aktionen}
  \dsaKartenstich{KaP-Kosten}{8 KaP}
  \dsaKartenstich{Reichweite}{Berührung}
  \dsaKartenfortsetzung
  \dsaKartenstich{Wirkungsdauer}{QS Spielrunden}
  \dsaKartenstich{Zielkategorie}{Zone}
  \dsaKartenstich{Merkmal}{Objekt}
  \dsaKartenstich{Verbreitung}{Geweihte des Flussvaters}
  \dsaKartenstich{Steigerungsfaktor}{C}
  \dsaKartentext{Wer die Zone betritt, legt eine Probe auf
    Körperbeherrschung ab, erschwert um QS. Misslingt sie, erhält er den
    Status \emph{Fixiert}, bis er sich befreit.}
\end{dsaKarteZauber}

% --- Sonderfertigkeit -------------------------------------------------

\begin{dsaKarteSonderfertigkeit}
  \dsaKartenname{Klingentanz}
  \dsaKartenuntertitel{Kampf-Sonderfertigkeit}
  \dsaKartennummer{UFR 09}

  \dsaKartentext{Wer den Klingentanz beherrscht, weicht einem Angriff aus
    und führt die eigene Waffe in derselben Bewegung weiter. Die Uferlande
    kennen den Kniff von den Flussgarden, die auf schwankenden Planken
    kämpfen lernen.}
  \dsaKartenstich{Regel}{Nach einer gelungenen Verteidigung mit einer
    Waffe der Kampftechnik Fechtwaffen darf sofort eine Attacke geführt
    werden, erschwert um 2. Dies gilt einmal je Kampfrunde.}
  \dsaKartenstich{Voraussetzungen}{GE 13, Fechtwaffen 10, Ausweichen 6}
  \dsaKartenstich{AP-Wert}{15 Abenteuerpunkte}
\end{dsaKarteSonderfertigkeit}

% --- Kultur -----------------------------------------------------------

\begin{dsaKarteKultur}
  \dsaKartenname{Uferlande (Stufe I)}
  \dsaKartenuntertitel{Wesenszug}
  \dsaKartennummer{UFR 10}

  \dsaKartenstich{Wasserfest}{Wer in den Uferlanden aufgewachsen ist,
    kennt das Wasser, bevor er laufen kann. Proben auf Schwimmen und
    Boote \& Schiffe sind um 1 erleichtert.}
  \dsaKartenstich{Fremdenscheu}{Die Uferländer halten zusammen und alle
    anderen auf Abstand. Die erste Probe auf Überreden oder Betören
    gegenüber einem Fremden ist um 1 erschwert.}
  \dsaKartenstich{Kombinierbar mit}{Stufe II: Moorkante, Stufe II:
    Wehrheimer Werft}
\end{dsaKarteKultur}

% ======================================================================
% Verbrauchsgegenstände: einseitig, mit Stückzahl — türkiser Rand
% ======================================================================
%
% Diese Karten haben keine eigene Rückseite. Sie bekommen den generischen
% Rücken aus der Klasse, damit ein Stapel gemischt werden kann.
%
% \dsaKartenanzahl gilt nur für den Druckbogen: dort erscheint die Karte so
% oft. In der Einzelkartenfassung bleibt es bei einer Seite je Kartenseite,
% denn das ist die Druckvorlage.

\dsaKartenfarbe{tuerkis}

\begin{dsaKarteVerbrauch}
  \dsaKartenanzahl{10}
  \dsaKartenname{Heiltrank}
  \dsaKartenuntertitel{Verbrauchsgegenstand, Stufe 1}
  \dsaKartenbild{platzhalter-heiltrank}
  \dsaKartennummer{UFR 13}

  \dsaKartentext{Heilt \dsaFormel{1W6} Lebenspunkte. Brauqualität 1,
    Preis 20 S.}
\end{dsaKarteVerbrauch}

\begin{dsaKarteVerbrauch}
  \dsaKartenanzahl{4}
  \dsaKartenname{Wirselkraut}
  \dsaKartenuntertitel{Verbrauchsgegenstand}
  \dsaKartenbild{platzhalter-wirselkraut}
  \dsaKartennummer{UFR 14}

  \dsaKartentext{Betäubt für \dsaFormel{2W6} Stunden. Eine Probe auf
    Selbstbeherrschung, erschwert um 2, halbiert die Dauer.}
\end{dsaKarteVerbrauch}

% ----------------------------------------------------------------------
% Die Rückenmotive zum Vergleich
% ----------------------------------------------------------------------
%
% Drei weitere einseitige Karten, nur damit jedes Rückenmotiv einmal im
% gebauten PDF steht. Die Rückseite ist immer die Seite NACH der Karte.

\dsaKartenfarbe{rot}

\begin{dsaKarteVerbrauch}
  \dsaKartenname{Uferbrand}
  \dsaKartenuntertitel{Rücken: Raute in Randfarbe}
  \dsaKartenbild{platzhalter-uferbrand}
  \dsaKartennummer{UFR 15}

  \dsaKartentext{Ein Schluck Kornbrand aus Wehrheim. Erleichtert die
    nächste Probe auf Willenskraft um 1 und erschwert alle übrigen für
    eine Stunde um ebenso viel.}
\end{dsaKarteVerbrauch}

\dsaKartenfarbe{violett}
\dsaKartenruecken{\dsaRueckenRautenfeld}

\begin{dsaKarteVerbrauch}
  \dsaKartenname{Schlafmohnkugel}
  \dsaKartenuntertitel{Rücken: Rautenfeld statt Raute}
  \dsaKartenbild{platzhalter-schlafmohnkugel}
  \dsaKartennummer{UFR 16}

  \dsaKartentext{Wer sie zerbeißt, schläft \dsaFormel{1W6} Stunden. Wecken
    gelingt nur mit einer Probe auf Heilkunde Gift, erschwert um 4.}
\end{dsaKarteVerbrauch}

% Dieselbe Raute, aber auf der bandlosen Fläche: das Schuppenband lässt
% sich nicht abschalten — es steckt mit allem anderen in einer flachen
% Datei —, aber werkzeuge/kartengrafik.py leitet eine Fassung ohne Band ab.
\dsaKartenfarbe{bernstein}
\dsaKartenruecken{\dsaRueckenRaute}
\dsaKartenrueckflaeche{spielkarte-ruecken}

\begin{dsaKarteVerbrauch}
  \dsaKartenname{Flussvaterwasser}
  \dsaKartenuntertitel{Rücken: Raute ohne Schuppenband}
  \dsaKartenbild{platzhalter-flussvaterwasser}
  \dsaKartennummer{UFR 17}

  \dsaKartentext{Geweihtes Wasser aus dem Tempel zu Wehrheim. Hebt einen
    Zustand der Stufe 1 auf, gleich welcher Art.}
\end{dsaKarteVerbrauch}

% Für die Karten danach wieder die Vorgabe: Raute auf der Kartenfläche.
\dsaKartenrueckflaeche{}

% ======================================================================
% Vierte Gruppe: was im Moor lauert — bernsteinfarbener Rand
% ======================================================================
%
% Ab hier wieder zweiseitige Karten, also zurück zum eigenen Rückseiteninhalt.

\dsaKartenfarbe{bernstein}

% --- Monster ----------------------------------------------------------

\begin{dsaKarteMonster}
  \dsaKartenname{Moorschrecke}
  \dsaKartenuntertitel{Tier, nicht humanoid}
  \dsaKartenbild{platzhalter-moorschrecke}
  \dsaKartennummer{UFR 11}
  \dsaKartenmedaillon{platzhalter-moorschrecke-rund}

  \dsaKartenattribut{MU}{11}   \dsaKartenattribut{KL}{3 (t)}
  \dsaKartenattribut{IN}{13}   \dsaKartenattribut{CH}{8}
  \dsaKartenattribut{FF}{6}    \dsaKartenattribut{GE}{15}
  \dsaKartenattribut{KO}{12}   \dsaKartenattribut{KK}{11}
  \dsaKartenattribut{LeP}{24}  \dsaKartenattribut{AsP}{\dsaMinus}
  \dsaKartenattribut{KaP}{\dsaMinus}
  \dsaKartenattribut{INI}{\dsaFormel{14+1W6}}
  \dsaKartenattribut{VW}{8}    \dsaKartenattribut{SK}{2}
  \dsaKartenattribut{ZK}{3}    \dsaKartenattribut{GS}{10}
  \dsaKartenattribut{AT}{13}   \dsaKartenattribut{PA}{9}
  \dsaKartenattribut{TP}{\dsaFormel{1W6+3}}
  \dsaKartenattribut{RW}{kurz}

  \dsaKartenstich{RS/BE}{2/0}
  \dsaKartenstich{Aktionen}{1}
  \dsaKartenstich{Vorteile/Nachteile}{Herausragender Sinn (Gehör),
    Dunkelsicht}
  \dsaKartenstich{Talente}{Körperbeherrschung 9, Schwimmen 12,
    Sinnesschärfe 10, Verbergen 14}
  \dsaKartenstich{Größenkategorie}{mittel}
  \dsaKartenstich{Sonderregeln}{\emph{Sprung aus dem Wasser} — im ersten
    Angriff einer Kampfrunde aus dem Wasser heraus ist die Verteidigung
    um 2 erschwert.}
\end{dsaKarteMonster}

\begin{dsaKarteMonster}
  \dsaKartenname{Nebelschleicher}
  \dsaKartenuntertitel{Chimäre, nicht humanoid}
  \dsaKartenbild{platzhalter-nebelschleicher}
  \dsaKartennummer{UFR 12}
  \dsaKartenmedaillon{platzhalter-nebelschleicher-rund}

  \dsaKartenattribut{MU}{16}   \dsaKartenattribut{KL}{9}
  \dsaKartenattribut{IN}{14}   \dsaKartenattribut{CH}{12}
  \dsaKartenattribut{FF}{10}   \dsaKartenattribut{GE}{14}
  \dsaKartenattribut{KO}{15}   \dsaKartenattribut{KK}{16}
  \dsaKartenattribut{LeP}{45}  \dsaKartenattribut{AsP}{20}
  \dsaKartenattribut{KaP}{\dsaMinus}
  \dsaKartenattribut{INI}{\dsaFormel{15+1W6}}
  \dsaKartenattribut{VW}{9}    \dsaKartenattribut{SK}{4}
  \dsaKartenattribut{ZK}{5}    \dsaKartenattribut{GS}{8}
  \dsaKartenattribut{AT}{15}   \dsaKartenattribut{PA}{11}
  \dsaKartenattribut{TP}{\dsaFormel{2W6+2}}
  \dsaKartenattribut{RW}{mittel}

  \dsaKartenstich{RS/BE}{4/0}
  \dsaKartenstich{Aktionen}{2}
  \dsaKartenstich{Vorteile/Nachteile}{Dunkelsicht, Immunität gegen Kälte}
  \dsaKartenstich{Talente}{Einschüchtern 12, Verbergen 16, Willenskraft 11}
  \dsaKartenstich{Größenkategorie}{groß}
  \dsaKartenstich{Sonderregeln}{\emph{Nebelleib} — außerhalb von Nebel
    verliert der Nebelschleicher 1~LeP je Spielrunde und kann nicht
    zaubern.}
\end{dsaKarteMonster}
